\section{Cancelación a todos los órdenes fijos necesarios en un bosque}
Consideremos primero sólo interacciones ordinarias. En el producto local inclinado
por $R'(m)$, sea $Z=X-m$ y $V_m$ la covarianza, diagonal por grupos. Para bosques,
el error se escribe como logaritmo de esperanza del exponencial de
$R(m+Z)-R(m)-R'(m)Z+A_*(m+Z)-A_*(m)$ más
$\tfrac12\log\det(I-V_mR''(m))$, salvo constantes comunes.
A campo inclinado congelado, cada término conexo de ese logaritmo tiene $p$ patas
centradas, $q$ factores de media, $b$ bloques de cumulantes y $N_R,N_A$ vértices de
acción y prefactor. La conexidad sigue de la inversión momento--cumulante; cada bloque
está en un único grupo. El grafo con una arista por pata centrada tiene
$L=p-N_R-N_A-b+1\ge0$ y $p+q\le w+N_R$. Luego su soporte es a lo sumo
$q+b\le w+1-N_A-L$. Los bloques de cumulantes y los vértices $R$ tienen grado
$\ge2$, y los $A_*$ grado $\ge1$. Si $N_A=1$ no puede haber un árbol con una sola
hoja posible. Si $N_A=0$, un único ciclo independiente fuerza todos los grados
a ser dos. Esos términos suman $-\log\det(I-V_mR''(m))/2$ por las contracciones gaussianas,
y se cancelan. Todo término ordinario superviviente tiene a lo sumo $w-1$ grupos.

Al reponer la silla, cada orden suplementario añade a lo sumo ese número de grupos.
En efecto, en la ecuación del campo de un grupo raíz $i$, el derivar un polinomio
$R_l$ deja a lo sumo $l$ factores de media. Las decoraciones de órdenes $t_a$ de cada
factor añaden, inductivamente, a lo sumo $\sum t_a$ grupos. El coste es $l+\sum t_a$.
Las derivadas de un cumulante local se toman en los campos de su propio grupo y sus
decoraciones comparten esa raíz. Esto prueba la afirmación por inducción para cada
orden considerado. Este argumento trata la parte exclusivamente ordinaria.

\paragraph{Con una interacción excepcional.}
La cuenta de la parte 3 da $d\le w-1$ si hay dos excepcionales (cada una de déficit
al menos uno) o una de déficit $\ge2$. Para una sola pareja excepcional, obtener
$d=w$ exige $L_I=0$, déficit exactamente uno en esa pareja y cero en cada otro vértice.
Estos últimos son interacciones ordinarias de acción, productos de $l+1$ formas lineales
sobre sendos grupos distintos. El multiplicador excepcional también es producto de
formas lineales en grupos distintos, ajenos a las dos fibras distinguidas, que están
en grupos diferentes. Todas las aseveraciones valen monomio por monomio por multilinealidad.

\paragraph{Eliminación de una rama.}
En ese árbol enraizado en la excepción, tomamos una interacción ordinaria más alejada.
Tiene un único grupo de enlace y los demás son grupos hoja donde todavía no hay otra
interacción. Es el producto de una forma lineal del grupo de enlace y una por cada hoja.
La derivada al anular su amplitud equivale al producto de las medias de las formas
de las hojas por la derivada respecto del campo lineal en el grupo de enlace.
Para el exacto se factoriza la esperanza antes de derivar en las amplitudes que restan.
Para el canónico, el valor estacionario da la misma sustitución por el teorema de
la envolvente. La respuesta de la silla en el núcleo es la de aquel campo: al anular
la amplitud, el jacobiano separa núcleo y hojas. En el prefactor, la derivada explícita
de la Hessiana del producto sólo tiene bloques cruzados, de traza nula contra la inversa
en bloques. En las hojas forestales la corrección local del determinante es idénticamente
uno en todo campo. Las derivadas implícitas en el núcleo son exactamente las inducidas
por el campo citado. En estas hojas la media de la silla es la verdadera.
Los campos pueden ser arbitrarios en el polidisco; al repetir la eliminación también
se derivan los factores de hojas ya obtenidos por la regla del producto.
Así el coeficiente buscado consiste en derivadas de la diferencia en la raíz y factores
de medias de hojas y sus derivadas, comunes a las dos evaluaciones.

En la raíz, a primer orden en la amplitud común, $P(X)+B(X)=0$ en átomos binarios.
Sea $m_i=\E y_i$, $v_i=m_i(1-m_i)$ la varianza real de la fibra categórica en el
bosque. Por independencia entre las dos fibras y por $F(m_i)=-m_i-v_i$,
\[
 P(m)+B(m)+\tfrac12\tr(V_m P''(m))
 =2F_iF_h+(m_i+v_i)F_h+(m_h+v_h)F_i=0 .
\]
Los otros factores lineales del multiplicador sólo aportan sus medias: sus grupos
son distintos y la traza no ve bloques cruzados. Todas las derivadas de esta identidad
en campos siguen siendo cero. Desaparece así el único caso excepcional de soporte $w$.
La primera contribución posible sobre $c$ grupos forestales tiene orden $s^{c+1}$.
